\section{Implementation and Experiments}
This section outlines the pipeline developed for the project, describing both the training phase and the C++ implementation, as well as some of the design choices.
Check out the doxygen documentation in \texttt{doxygen.txt} and the source code for more details.

\subsection{Training of the model}
In this project, the YOLO11s model was selected, to deliver good performance with enough speed.
The YOLO model was trained on the dataset using two RTX 3090 GPUs inside the UniPD DEI cluster.
The training dataset was split into 3 smaller subsets such that:
\begin{itemize}
    \item Each subset contains images (and corresponding labels) which are different than the other datasets (preventing data leakage);
    \item Each subset has a maximum size which allows each one of them to be loaded into RAM.
\end{itemize}
As a result, the training process was divided into three sequential steps: each step loaded the model weights from the 
previous step and continued training on the next subset. This strategy allowed for a faster and more fault-tolerant training process 
with a negligible penalty on accuracy. The full script is available at \texttt{scripts/yolo\_training.py}. \\ \newline
The model was specifically trained to detect bounding boxes enclosing the suit and rank on the corners of a 
poker card. This means that at most 2 bounding boxes can be found for the same card, ensuring that even if a card is partially occluded, 
the system remains robust as long as at least one corner is visible. \\ \newline
Since the final application is developed in C++, the model had to be exported into a ONNX format and imported 
into the program using a library like ONNX Runtime.

\subsection{Inference pipeline}
The application is organized into modular components within the \texttt{src/} directory. The execution flow 
is coordinated by \texttt{finalProject.cpp}, which acts as the main entry point, allowing the user to choose between camera 
input or file processing. The core pipeline is implemented as follows:
\begin{enumerate}
    \item Input Handling and Pre-processing: user interaction and terminal input parsing are handled in \texttt{hermann\_serain.cpp}. 
    Once the source is selected, \texttt{marco\_annunziata.cpp} prepares the image. The image is formatted as a valid input of the 
    Yolo11s model and the ONNX Runtime session. Letterbox padding of the input images is enabled by default, since YOLO models use 
    it for the training phase, but the option can also be disabled when calling the function used for the inference;
    \item Model Inference: the \texttt{YOLO\_model} class in \texttt{marco\_annunziata.cpp} manages the ONNX Runtime session. 
    It handles the inference execution, which generates approximately 84,000 potential detections;
    \item Post-processing: each detection result consists of 4 values defining the object's bounding box (position and size), and 
    a confidence score for every possible class which the object might correspond to. The system uses Non-Maxima Suppression in 
    \texttt{marco\_annunziata.cpp} to keep only 1 bounding box for each detected object.
\end{enumerate}
In order to handle detection of small cards in big images, a sliding window approach is used: the image is subdivided into smaller, slightly 
overlapping tiles, with size equal to the image size used for training the model (640px). The overhead introduced by this method is mitigated 
by using multi-threading, where each tile is processed by a separate thread in parallel.

\subsection{Hi-Lo Logic and Visualization}
After detection, the application logic is managed by \texttt{sveva\_turola.cpp}. This phase involves:
\begin{itemize}
    \item Card Grouping: the \texttt{cardValues} function groups nearby corner detections (suit and rank) belonging to the same 
    card into a single, unified bounding box. This grouping is governed by a dynamic distance threshold calculated based on the 
    symbol's diagonal, ensuring the system adapts to different card sizes and distances from the camera;
    \item Hi-Lo Counting System: bounding boxes are classified and color-coded: Low (2-6) in Green, Neutral (7-9) in 
    Blue, and High (10-A) in Red. The system maintains a Running Count displayed in real-time on the GUI.
\end{itemize}
The module handles both live camera feeds and video files, saving processed frames only when the detection count changes. 
The final result is rendered in a fullscreen window with custom playback controls (\texttt{p} to pause, \texttt{q} to quit).

\subsection{Performance metrics}
A complex but fast evaluation system is integrated in \texttt{hermann\_serain.cpp}. The program will compare the predictions made by the model against the ground truth labels, only if ground truths are provided. The default image and video bundled inside the project are already labeled,
but it's also possible to add labels for custom images and videos. For more information, see the README.md file.
\\\\
The pipeline to compute the metrics relies on counting the number of True Positives (TP), False Positives (FP) and False Negatives (FN) over all card detections (there are no True Negatives since everything that is not a card is a TN).
One TP is considered as a full card detection, but the model detects the rank and suit of all corners of the card to improve robustness to occlusions, meaning that the system has to be able to aggregate the detections in the separate corners of each card into a single detection of the full card.
\\\\
Another increase in complexity comes from the problem that videos need a labeling format, the MOT format, which is different than the YOLO format used for images.
There are multiple reasons for why the MOT format has been chosen for videos:
\begin{itemize}
    \item It's a format compatible with the video labeling software used (CVAT), particularly with the visiblity feature that keeps note of occlusions during the tracking of the bounding boxes;
    \item Each video corresponds in this format to a single label file, which is more compact and convenient than using one file per video frame;
    \item The format of the detections in the file are easily readable line-by-line, in contrast with the XML formats that could be harder to parse correctly for the amount of information that is actually needed.
\end{itemize}
For each frame of the video, the system:
\begin{itemize}
    \item Discards the ground truths where the visibility of the bounding box is 0;
    \item Checks if the bounding boxes are in the same position of the ground truths, and sees if they are of the same class. An overlap between boxes is detected if their Intersection over Union (IoU) is over a threshold of 0.5;
    \item Aggregates the bounding boxes corresponding to corners of the same card to a single full card detection;
    \item Updates the counters for TPs, FPs, FNs.
\end{itemize}
The IoU is the ratio between the intersection of the predicted and actual bounding boxes and the total area covered by both boxes:
\[ IoU = \frac{\text{Area of Intersection}}{\text{Area of Union}} \]
The conunters of TPs, FPs and FNs for each image/video are finally used to compute the metrics. The following metrics quantitatively evaluate the performance of the system:
\begin{itemize}
    \item Precision: measures the proportion of correctly predicted bounding boxes among all detections;
    \[ Precision = \frac{TP}{TP + FP} \]
    \item Recall: measures the proportion of correctly predicted bounding boxes among all ground truth boxes;
    \[ Recall = \frac{TP}{TP + FN} \]
    \item F1-Score: harmonic mean of Precision and Recall used to evaluate how well a classification model performs.
    \[  F1 = 2 \cdot \frac{\text{Precision} \cdot \text{Recall}}{\text{Precision} + \text{Recall}} \]
\end{itemize}
The final performance statistics are shown once the video or image processing is complete.